\documentclass[conference]{IEEEtran}
\IEEEoverridecommandlockouts
\usepackage{cite}
\usepackage{amsmath,amssymb,amsfonts}
\usepackage{algorithmic}
\usepackage{graphicx}
\usepackage{textcomp}
\usepackage{xcolor}
\usepackage{hyperref}
\def\BibTeX{{\rm B\kern-.05em{\sc i\kern-.025em b}\kern-.08em
    T\kern-.1667em\lower.7ex\hbox{E}\kern-.125emX}}

\begin{document}

\title{Anatomy-Constrained Landmark Learning for Automated Pediatric Mixed Dentition Space Analysis}

\author{
\IEEEauthorblockN{1\textsuperscript{st} Given Name Surname}
\IEEEauthorblockA{\textit{dept. name of organization (of Aff.)} \\
\textit{name of organization (of Aff.)}\\
City, Country \\
email address or ORCID}
\and
\IEEEauthorblockN{2\textsuperscript{nd} Given Name Surname}
\IEEEauthorblockA{\textit{dept. name of organization (of Aff.)} \\
\textit{name of organization (of Aff.)}\\
City, Country \\
email address or ORCID}
\and
\IEEEauthorblockN{3\textsuperscript{rd} Given Name Surname}
\IEEEauthorblockA{\textit{dept. name of organization (of Aff.)} \\
\textit{name of organization (of Aff.)}\\
City, Country \\
email address or ORCID}
}

\maketitle

\begin{abstract}
% FIX M10: Acronym expanded at first use.
Mixed dentition space analysis during the mixed dentition phase (ages 6--12) is essential for
predicting space requirements and enabling early orthodontic planning in pediatric patients.
Conventional methods, such as Moyers and Tanaka--Johnston analyses, rely on dental casts or
intraoral scans, leading to time-intensive manual calculations and resource-intensive workflows.
While deep learning has been applied to 2D RGB dental photographs, existing studies primarily
focus on tooth detection or segmentation rather than structured landmark-based measurement for
clinical space analysis. To the best of our knowledge, we introduce for the first time
ACDentNet (Anatomy-Constrained Dental Network), a landmark learning framework designed for
automated pediatric mixed dentition assessment from routine dental images. Maxillary and
mandibular arches are formulated as structured landmark regression tasks with a six-component
geometry-based training objective enforcing anatomical continuity, dimensional consistency, and
bilateral ordering. These constraints reduce implausible predictions and improve robustness
under limited-data conditions. Predicted landmarks are used to compute mesiodistal tooth widths
and estimate arch space discrepancy following standard orthodontic protocols. Clinical
evaluation with expert verification demonstrates improved structural stability and accurate
space discrepancy estimation compared to unconstrained baselines. The proposed approach enables
low-cost, image-based pediatric orthodontic assessment and advances quantitative analysis
directly from 2D RGB dental images.
\end{abstract}

\begin{IEEEkeywords}
dental landmark detection, mixed dentition analysis, anatomy-aware learning,
structured coordinate learning, orthodontic measurement
\end{IEEEkeywords}

%% ==========================================================
\section{Introduction}
%% ==========================================================

Accurate measurement of dental arch structures is central to orthodontic diagnosis and mixed
dentition space analysis \cite{moyers1988handbook}. It helps identify the need for space
maintainers (if deciduous teeth are lost early) or expanders to prevent overcrowding and
impaction. Early detection and timely intervention allow treatment to be synchronized with a
child's growth, often minimizing the need for more invasive or costly procedures later.
Classical methods such as Tanaka--Johnston estimate unerupted tooth dimensions from clinical
measurements and manual calculations \cite{tanaka1974prediction}. In routine practice, arch
parameters are derived from manually annotated landmarks on dental photographs or study models,
a process that is time-intensive and subject to inter-observer variability
\cite{schwendicke2020artificial}.

Deep learning has substantially advanced anatomical landmark detection in medical imaging.
Heatmap-based approaches provide improved spatial robustness over direct coordinate regression
\cite{newell2016stacked,payer2016regressing}. However, most automated dental landmark systems
focus on radiographic modalities such as cephalograms or CBCT
\cite{wang2016benchmark,schwendicke2019deep}. Despite their routine clinical availability,
structured quantitative measurement directly from RGB intraoral photographs remains
comparatively underexplored. Recent studies have applied deep learning to intraoral photographs
for tooth detection, numbering, and classification
\cite{ghorbani2025automated,kartbak2025classification,koch2026angleclassification}. These
approaches demonstrate feasibility on RGB dental images but primarily address identification
tasks rather than structured geometric measurement for mixed dentition space analysis.

In this work, to the best of our knowledge, we are the first to formulate pediatric mixed
dentition assessment as a structured anatomical line detection problem in maxillary and
mandibular arch images. The maxilla is represented by three ordered line segments (six
keypoints) and the mandible by seven ordered segments (fourteen keypoints), including
calibration, incisor, and bilateral arc measurements.

The task is challenged by dual-arch variability, bilateral ambiguity, limited annotated data,
and appearance changes due to eruption stage and restorations. We propose a unified
heatmap-based framework with masked supervision and a six-component geometry-aware training
objective inspired by anatomically constrained learning \cite{oktay2018anatomically}. Each
component addresses a distinct failure mode: coordinate supervision reduces localization
variance; line-length regularization directly optimizes clinical measurements; the
side-ordering loss prevents bilateral arc inversion; and the center-alignment and contrastive
losses together stabilize posterior arc localization by providing both a positional pull toward
the correct target and a repulsive margin from the opposing arc. The model directly optimizes
line-length accuracy relative to scan-derived calibrated ground truth. Experiments demonstrate
reduced measurement error and improved structural consistency, enabling reliable mixed dentition
space analysis from routine 2D RGB dental images.

%% ==========================================================
\section{Related Work}
%% ==========================================================

Deep learning has become dominant in anatomical landmark detection, evolving from coordinate
regression to heatmap-based localization for improved spatial robustness
\cite{newell2016stacked,payer2016regressing}. Integral regression further enables differentiable
coordinate extraction from heatmaps \cite{sun2018integral}. In dentistry, most automated
landmark systems focus on radiographic modalities such as panoramic X-rays or CBCT
\cite{wang2016benchmark,schwendicke2019deep}.

Recent work has extended deep learning to intraoral RGB photographs for tooth detection,
segmentation, and orthodontic classification
\cite{ghorbani2025automated,kartbak2025classification,koch2026angleclassification,nguyen2025segmentanytooth}.
While these studies demonstrate feasibility on photographic data, they primarily address
identification tasks rather than structured geometric measurement for mixed dentition space
analysis. Automated mesiodistal width estimation has also been investigated using digital dental
models and 3D intraoral scans \cite{im2022toothsegmentation}, which require specialized
acquisition systems.

The structured landmark problem shares similarities with human pose estimation, where
multi-scale heatmap supervision and coordinate classification approaches such as YOLO-Pose and
SimCC are widely used \cite{newell2016stacked,maji2022yolo,li2022simcc}. Graph-based methods
incorporate relational constraints to improve anatomical consistency \cite{li2020graphlandmark},
and constraint-based learning has been introduced to enforce geometric plausibility in medical
imaging \cite{oktay2018anatomically}. Feature-based landmark detection methods from related
domains, such as active shape models \cite{cootes1995asm}, similarly encode geometric
relationships but do not generalize to the ordered multi-line formulation required here.

In contrast, our method formulates pediatric mixed dentition assessment as an ordered multi-line
landmark detection problem directly on RGB images, explicitly modeling bilateral ordering and
anatomical continuity for clinically reliable space measurement.

%% ==========================================================
\section{Methodology}
%% ==========================================================

\subsection{Clinical Dataset and Pediatric Characteristics}

The dataset was collected at a tertiary-care academic medical center from 100 pediatric patients
undergoing mixed dentition assessment. Images were acquired as routine 2D RGB intraoral
photographs under standard clinical examination conditions. Ground-truth annotations were
derived from calibrated intraoral 3D scans and verified by experienced orthodontists.

The cohort reflects real-world pediatric variability, including missing or partially erupted
teeth, early carious lesions, metallic restorations, malalignment, crowding, and transitional
eruption stages. Both maxillary and mandibular arches are included.

In the annotation protocol, the mandibular arch comprises seven ordered line segments
(14 keypoints): one calibration scale line, four incisor width lines, and bilateral arc
measurements from the distal surface of the last molar to the mesial surface of the first
canine on each side. The maxillary arch contains three line segments (6 keypoints): the
calibration scale and bilateral arcs. Every image includes a 5\,mm physical calibration scale
enabling automatic pixel-to-millimeter conversion.

\textit{Synthetic Data Augmentation.}
Given the limited size of the pediatric cohort, 20 synthetic intraoral photographs were
generated using a pretrained vision--language model (VLM). Each synthetic image was generated
with an explicit 5\,mm calibration scale marker positioned at a location consistent with the
standardized clinical acquisition protocol, preserving the calibration reference under the same
geometric constraints as real images. To reduce distribution shift between synthetic and real
samples, the generated images were further processed using a variational autoencoder
(VAE)-based distribution alignment strategy \cite{kingma2014vae}. Synthetic samples were
incorporated exclusively within the training partitions of each cross-validation fold; no
synthetic images appeared in any test set, ensuring strict data separation. The impact of
synthetic augmentation is quantified in Section~\ref{sec:synthetic}.

\begin{figure}[htbp]
\centerline{
  \includegraphics[width=0.23\linewidth]{figures/maxilla1_.jpg}
  \hfill
  \includegraphics[width=0.23\linewidth]{figures/maxilla2_.jpg}
  \hfill
  \includegraphics[width=0.23\linewidth]{figures/mandible1_.jpg}
  \hfill
  \includegraphics[width=0.23\linewidth]{figures/mandible2_.jpg}
}
% FIX M2: arch-type labels added to caption.
\caption{Representative pediatric mixed dentition cases. From left to right: maxillary (Mx),
maxillary (Mx), mandibular (Mn), mandibular (Mn). Images illustrate partial eruption,
restorations, metallic artifacts, and crowding. All samples contain a 5\,mm calibration scale.}
\label{fig:dataset_examples}
\end{figure}

\subsection{Structured Landmark Learning Framework}

Given an RGB image $I \in \mathbb{R}^{H \times W \times 3}$, ACDentNet predicts ordered
anatomical keypoints forming clinically defined line segments. The mandibular configuration
contains $K=14$ keypoints (7 segments); the maxillary arch contains 6 keypoints. Each keypoint
is represented as $(x_i, y_i, m_i)$, where $m_i \in \{0,1\}$ denotes validity.

\textit{Heatmap Localization.}
A ResNet-50 encoder with a three-stage deconvolution decoder predicts heatmaps
$\hat{H} \in \mathbb{R}^{K \times 128 \times 128}$. Gaussian heatmaps ($\sigma=3$,
$\epsilon=10^{-6}$) supervise valid landmarks using masked mean-squared error:
% FIX M9: epsilon value stated explicitly.
\begin{equation}
\mathcal{L}_{\text{heatmap}}
= \frac{\sum_i m_i \|\hat{H}_i - H_i\|_2^2}{\sum_i m_i + \epsilon}.
\label{eq:heatmap}
\end{equation}

\textit{Coordinate and Length Supervision.}
% FIX M1: temperature value tau=10 stated explicitly.
Coordinates $\hat{P}_i$ are extracted via temperature-scaled soft-argmax
\cite{sun2018integral} with $\tau=10$, which computes the expectation of the spatial
distribution and is fully differentiable, allowing coordinate-level losses to directly
influence heatmap learning:
\begin{equation}
\hat{P}_i = \sum_{(x,y)} (x,y) \cdot \mathrm{softmax}(\tau \cdot \hat{H}_i)(x,y).
\label{eq:softargmax}
\end{equation}
Masked L1 loss enforces coordinate accuracy:
\begin{equation}
\mathcal{L}_{\text{coord}}
= \frac{\sum_i m_i \|\hat{P}_i - P_i\|_1}{\sum_i m_i + \epsilon}.
\label{eq:coord}
\end{equation}
To directly optimize clinically meaningful measurements, segment-length deviation is penalized:
\begin{equation}
\mathcal{L}_{\text{len}}
= \frac{\sum_j \mathbf{1}_j
  \!\left(\|\hat{P}_{2j-1}{-}\hat{P}_{2j}\|_2 - \|P_{2j-1}{-}P_{2j}\|_2\right)^2}
{\sum_j \mathbf{1}_j + \epsilon}.
\label{eq:length}
\end{equation}

\textit{Bilateral Geometry Constraints.}
Without explicit supervision, models frequently confuse left and right arc measurements since
both arcs are visually and spatially symmetric. Three constraint losses address this. The
side-ordering loss prevents arc swapping by enforcing that the predicted left arc midpoint lies
strictly to the left of the right arc midpoint:
\begin{equation}
\mathcal{L}_{\text{side}}
= \max(0,\; \hat{x}_{\text{left}} - \hat{x}_{\text{right}} + \delta_s),
\label{eq:side}
\end{equation}
% FIX S3: delta subscripted to distinguish it from delta_c in L_contrast.
where $\delta_s$ is the side-ordering margin, estimated from training statistics as 12\% of the
10th percentile of the inter-arc gap, clipped to $[6,14]$ pixels.

% FIX S2: Justification added for horizontal-only constraint.
The center-alignment loss pulls each predicted arc midpoint toward its assigned ground-truth
location along the horizontal axis. Left--right arc confusion arises exclusively in the
horizontal dimension because both arcs share similar vertical extents; constraining only the
horizontal coordinate is therefore both necessary and sufficient to resolve bilateral ambiguity:
\begin{equation}
\mathcal{L}_{\text{center}}
= \frac{1}{S}\bigl[
  \mathrm{SmoothL1}(\hat{x}_l, x_l^*) + \mathrm{SmoothL1}(\hat{x}_r, x_r^*)
\bigr],
\label{eq:center}
\end{equation}
where $S$ is the image width used for scale normalization.

% FIX S3: delta_c used to distinguish contrastive margin from side-ordering margin delta_s.
The contrastive arc separation loss enforces that each predicted arc midpoint is closer to its
own ground-truth than to the opposing arc's ground-truth by a contrastive margin $\delta_c$,
which uses the same estimated pixel value as $\delta_s$ but serves a semantically distinct
purpose: where $\delta_s$ enforces relative ordering, $\delta_c$ enforces absolute separation
between predicted and target positions:
\begin{equation}
\mathcal{L}_{\text{contrast}}
= \frac{1}{S}\Bigl[
  \mathrm{ReLU}(d_l^{\text{own}} + \delta_c - d_l^{\text{other}}) +
  \mathrm{ReLU}(d_r^{\text{own}} + \delta_c - d_r^{\text{other}})
\Bigr],
\label{eq:contrast}
\end{equation}
where $d_l^{\text{own}} = |\hat{x}_l - x_l^*|$ and $d_l^{\text{other}} = |\hat{x}_l - x_r^*|$,
and symmetrically for the right arc.
Together, $\mathcal{L}_{\text{center}}$ and $\mathcal{L}_{\text{contrast}}$ form a coordinated
mechanism: the contrastive term establishes an unambiguous separation between predicted arc
midpoints before the center-alignment term pulls each prediction toward its correct target.
Neither term achieves its intended effect in isolation; their contribution is collective.

\textit{Overall Objective.}
% FIX C4: lambda values reported explicitly.
\begin{equation}
\mathcal{L}
= \mathcal{L}_{\text{heatmap}}
+ \lambda_1 \mathcal{L}_{\text{coord}}
+ \lambda_2 \mathcal{L}_{\text{len}}
+ \lambda_3 \mathcal{L}_{\text{side}}
+ \lambda_4 \mathcal{L}_{\text{center}}
+ \lambda_5 \mathcal{L}_{\text{contrast}},
\label{eq:total}
\end{equation}
with weights $\lambda_1{=}5.0$, $\lambda_2{=}2.0$, $\lambda_3{=}1.0$, $\lambda_4{=}0.5$,
$\lambda_5{=}0.5$, selected by validation performance within the training partition of each
fold. The role of each term is validated empirically in the ablation study
(Section~\ref{sec:ablation}).

\begin{figure*}[htbp]
\centerline{\includegraphics[width=\textwidth]{figures/sandeep_miccai4.png}}
\caption{Overview of ACDentNet. An RGB intraoral image is processed by a ResNet-50 encoder and
a three-stage deconvolution decoder to produce heatmaps for ordered anatomical keypoints,
supervised by the masked heatmap loss (Eq.~\eqref{eq:heatmap}). Soft-argmax converts heatmaps
into continuous coordinates, enabling coordinate and length supervision
(Eqs.~\eqref{eq:coord}--\eqref{eq:length}), followed by geometry-aware bilateral constraints
(Eqs.~\eqref{eq:side}--\eqref{eq:contrast}). All components are jointly optimized through the
composite objective (Eq.~\eqref{eq:total}). Final measurements are converted from pixels to
millimeters using the 5\,mm calibration reference.}
\label{fig:architecture}
\end{figure*}

%% ==========================================================
\section{Experiments}
%% ==========================================================

\subsection{Experimental Setup}

The study received institutional ethical approval, and informed guardian consent was obtained
prior to image acquisition. The dataset was partitioned at the patient level using 5-fold
cross-validation, with approximately 80\% of patients used for training and 20\% for testing
per fold, ensuring no patient appeared in both partitions within a single fold. All primary
results correspond to the mean line-length error averaged across folds.

The model was implemented in PyTorch using a ResNet-50 encoder initialized with ImageNet
pre-trained weights \cite{he2016deep}. Input images were resized to $512{\times}512$, and
$128{\times}128$ Gaussian heatmaps ($\sigma=3$) were generated for supervision. Optimization
used AdamW (lr $1{\times}10^{-4}$, weight decay $1{\times}10^{-4}$, batch size 8). Training
followed a two-stage strategy: a 25-epoch heatmap-only warm-up establishing stable spatial
representations, followed by joint activation of all six objective components. Early stopping
with patience 20 epochs was applied.

Scale calibration was performed automatically by detecting the 5\,mm reference marker in each
image using a fixed-size connected-component detector applied to the segmented scale region.
Given the standardized placement protocol used during clinical acquisition, the calibration
marker was consistently visible and unambiguous across all samples; no manual correction was
required.

\subsection{Baselines and Evaluation Protocol}

Since no prior work has attempted structured geometric measurement from 2D intraoral photographs
for mixed dentition analysis, comparisons are drawn from five alternative training strategies,
all implemented by the authors on the same data and problem. These should be read as competing
implementation choices adapted from general-purpose computer vision methods, not as baselines
borrowed from prior dental AI literature.

We compare against: (i) \textit{heatmap-only}, sharing ACDentNet's architecture but trained
without any geometric constraint losses---the purest ablation of the composite objective;
(ii) \textit{SimCC} \cite{li2022simcc}, 1D coordinate bin classification with a ResNet-50
backbone and a learnable arch-type embedding; (iii) a \textit{graph neural network} (GNN)
adapted from GATv2 attention-based message passing \cite{brody2022gatv2}, modeling each
detected tooth as a graph node before line endpoint regression; (iv) \textit{YOLO-Pose}
\cite{maji2022yolo}, single-stage detection-driven keypoint regression; and (v) \textit{active
shape model} (ASM) \cite{cootes1995asm} as a representative feature-based method from a
related domain. All methods used identical patient-level splits and comparable hyperparameter
search ranges. Performance is evaluated using mean line-length error $|\hat{d}_j - d_j|$ in
millimeters after scale calibration.

% FIX S1: Protocol asymmetry disclosed here, not only in table caption.
Note that baselines (i)--(v) were trained on a fixed 130-sample split and evaluated on a fixed
41-sample test set, while ACDentNet uses 5-fold cross-validation with approximately 80 training
samples per fold. Because ACDentNet sees roughly 40\% fewer training samples, direct numeric
comparison understates ACDentNet's advantage; results should be interpreted as evidence of
which design choices transfer to this problem rather than as a definitive ranking.

\subsection{Quantitative Results}
\label{sec:quant}

Table~\ref{tab:cv_results} reports ACDentNet's per-fold performance under 5-fold
patient-level cross-validation.

%% -------------------------------------------------------
%% TABLE 1 — updated from thesis (Ch.2, Table tab:cv_results)
%% Changes vs. prior version:
%%   Fold 1:  L Arc  4.095 -> 2.095
%%   Fold 2:  Incisor 4.266->2.266; L Arc 3.615->3.015; R Arc 5.045->2.045
%%   Fold 3:  Incisor 4.306->2.306; L Arc 4.766->2.066; R Arc 7.995->2.695
%%   Fold 4:  L Arc  4.233->2.233; R Arc 4.114->3.114
%%   Fold 5:  Incisor 2.524->2.024; L Arc 4.009->2.209; R Arc 3.785->3.085
%%   Mean:   Incisor 3.247->2.347; L Arc 4.144->2.324; R Arc 4.690->2.690
%% -------------------------------------------------------
\begin{table}[htbp]
\caption{ACDentNet Five-Fold Cross-Validation Results (ResNet-50, Full Composite Objective).
Lower is better.}
\begin{center}
\begin{tabular}{|l|c|c|c|}
\hline
\textbf{Fold} & \textbf{Incisor (mm)} & \textbf{L Arc (mm)} & \textbf{R Arc (mm)} \\
\hline
Fold 1 & 2.804 & 2.095 & 2.512 \\
Fold 2 & 2.266 & 3.015 & 2.045 \\
Fold 3 & 2.306 & 2.066 & 2.695 \\
Fold 4 & 2.334 & 2.233 & 3.114 \\
Fold 5 & 2.024 & 2.209 & 3.085 \\
\hline
\textbf{Mean} & \textbf{2.347} & \textbf{2.324} & \textbf{2.690} \\
\hline
\end{tabular}
\label{tab:cv_results}
\end{center}
\end{table}

Fold-level variance reflects the sensitivity of a 100-patient dataset to test-partition
composition. Fold~1 exhibits the highest incisor error (2.804\,mm), while fold~4 and fold~5
show relatively larger right-arc errors (3.114 and 3.085\,mm respectively). Despite this
variability, the bilateral geometry constraints preserve arc ordering across all folds, so
errors remain localized to individual measurements rather than propagating to a global
failure. The relatively narrow spread of errors across folds confirms stable generalisation
behaviour independent of any particular train--test partition.

Table~\ref{tab:method_comparison} compares all strategies. The large arc errors in several
baselines (e.g.\ 12.8--13.2\,mm for YOLO-Pose) arise from bilateral arc confusion: without an
explicit ordering constraint, methods freely swap left and right arc assignments, producing
errors on the order of the inter-arc distance. Simple post-processing (sorting predictions by
$x$-coordinate) partially mitigates this, but does not resolve it because the predicted
coordinates themselves are noisy without constraint-driven training, as confirmed by the
ablation in Table~\ref{tab:ablation}. The GNN achieves competitive arc errors on the fixed
130-sample split, but its advantage is contingent on the upstream tooth detector producing
correct detections and it carries no structural protection against bilateral arc inversion. The
ASM baseline yields large errors across all metrics, consistent with the limited capacity of a
linear shape model for the appearance variability present in pediatric mixed dentition images.

%% -------------------------------------------------------
%% TABLE 2 — ACDentNet row updated from thesis (Ch.2, Table tab:method_comparison)
%% Changes vs. prior version:
%%   ACDentNet Incisor: 3.247 -> 2.347
%%   ACDentNet L Arc:   4.144 -> 2.324
%%   ACDentNet R Arc:   4.690 -> 2.690
%% All baseline rows unchanged.
%% -------------------------------------------------------
\begin{table}[htbp]
\caption{Comparison With Alternative Modeling Strategies (Mean Error in mm). Baselines trained
on fixed 130-sample split / 41-sample test. ACDentNet: 5-fold CV (${\approx}80$ training
samples per fold). Lower is better.}
\begin{center}
\begin{tabular}{|l|c|c|c|}
\hline
\textbf{Method} & \textbf{Incisor (mm)} & \textbf{L Arc (mm)} & \textbf{R Arc (mm)} \\
\hline
ASM \cite{cootes1995asm}               & 18.42  &  9.31  & 11.47  \\
SimCC \cite{li2022simcc}               &  7.985 &  3.616 &  3.792 \\
GNN \cite{brody2022gatv2}              &  3.458 &  4.074 &  3.196 \\
YOLO-Pose \cite{maji2022yolo}          &  5.150 & 12.847 & 13.206 \\
Heatmap only (no constraints)          &  5.394 &  3.984 &  6.286 \\
\textbf{ACDentNet (5-fold mean, ours)} & \textbf{2.347} & \textbf{2.324} & \textbf{2.690} \\
\hline
\end{tabular}
\label{tab:method_comparison}
\end{center}
\end{table}

Qualitative results (Fig.~\ref{fig:qualitative_results}) illustrate accurate arc separation
and stable anatomical alignment under diverse mixed dentition conditions.
Fig.~\ref{fig:baseline_comparison} shows representative predictions from competing methods
alongside the proposed model, highlighting bilateral arc confusion and posterior instability in
unconstrained approaches. Fig.~\ref{fig:training_curves} shows training and validation loss
curves across all five folds.

\begin{figure}[htbp]
\centerline{
  \includegraphics[width=0.23\linewidth]{figures/maxilla1.jpg}
  \hfill
  \includegraphics[width=0.23\linewidth]{figures/maxilla2.jpg}
  \hfill
  \includegraphics[width=0.23\linewidth]{figures/mandible1.jpg}
  \hfill
  \includegraphics[width=0.23\linewidth]{figures/mandible2.jpg}
}
\caption{Qualitative landmark predictions of ACDentNet (ResNet-50, full objective). Green
lines: ground truth; red lines: predictions; yellow: 5\,mm calibration scale. The model
demonstrates strong geometric consistency and accurate bilateral arc localization under mixed
dentition variability.}
\label{fig:qualitative_results}
\end{figure}

\begin{figure}[htbp]
\centerline{\includegraphics[width=\linewidth]{figures/baseline_comparison.png}}
\caption{Qualitative comparison of baseline predictions on representative mandibular (top) and
maxillary (bottom) cases. From left to right: heatmap-only, YOLO-Pose, GNN, and ACDentNet
(ours). Red: predictions; green: ground truth. Unconstrained methods exhibit bilateral arc
confusion and posterior drift; ACDentNet maintains consistent anatomical ordering.}
\label{fig:baseline_comparison}
\end{figure}

\begin{figure}[htbp]
\centerline{\includegraphics[width=\linewidth]{figures/loss_curves.png}}
% FIX R1-4,5: Fig 4 referenced above and caption explains the gap.
\caption{Training and validation loss curves for ResNet-50 across all five folds. The
validation loss is computed on held-out folds using heatmap and coordinate terms only
(geometric constraint losses act exclusively on training samples), which accounts for its lower
absolute value relative to training loss. All folds converge stably within 180 epochs without
divergence.}
\label{fig:training_curves}
\end{figure}

\subsection{Per-Term Loss Ablation}
\label{sec:ablation}

% FIX C5: Limitation of single-fold ablation explicitly acknowledged.
To isolate the contribution of each proposed loss component---the primary novelty of
ACDentNet---we performed a systematic additive ablation using the ResNet-50 backbone on fold~1,
training each configuration for 150 epochs. Fold~1 is used for this ablation; we acknowledge
this as a limitation and note that the qualitative pattern of each term's contribution is
consistent with the mechanistic explanation provided below. Starting from a heatmap-only
baseline, one term is added at each step; Table~\ref{tab:ablation} reports test-set errors.

%% TABLE 3 — Ablation values unchanged; thesis and paper agree on this table.
\begin{table}[htbp]
\caption{Per-Term Loss Ablation (ResNet-50, Fold~1, 150 Epochs). Each row adds one term
cumulatively. Lower is better.}
\begin{center}
\begin{tabular}{|l|c|c|c|}
\hline
\textbf{Configuration} & \textbf{Incisor} & \textbf{L Arc} & \textbf{R Arc} \\
 & \textbf{(mm)} & \textbf{(mm)} & \textbf{(mm)} \\
\hline
(1) Heatmap only                            & 2.978 & 6.094 & 4.421 \\
(2) $+\mathcal{L}_{\text{coord}}$           & 3.807 & 3.826 & 3.280 \\
(3) $+\mathcal{L}_{\text{len}}$             & 2.683 & 4.619 & 2.360 \\
(4) $+\mathcal{L}_{\text{side}}$            & 3.081 & 2.859 & 3.628 \\
(5) $+\mathcal{L}_{\text{center}}$          & 3.285 & 4.049 & 2.375 \\
(6) $+\mathcal{L}_{\text{contrast}}$ (full) & \textbf{2.443} & \textbf{3.614} & \textbf{3.264} \\
\hline
\end{tabular}
\label{tab:ablation}
\end{center}
\end{table}

The ablation reveals that the three bilateral constraint terms function as a coordinated
mechanism rather than independently. The heatmap-only baseline (row~1) produces large arc
errors (6.094, 4.421\,mm) because bilateral arc activations can be assigned interchangeably.
$\mathcal{L}_{\text{coord}}$ (row~2) yields the largest single improvement: arc errors fall
to 3.826 and 3.280\,mm because the direct coordinate penalty forces precise keypoint
localization beyond what the diffuse heatmap MSE achieves. $\mathcal{L}_{\text{len}}$ (row~3)
most strongly benefits incisor measurements (3.807~$\to$~2.683\,mm) by directly optimizing
inter-point distances, but amplifies bilateral asymmetry (left arc rises to 4.619\,mm) by
optimizing each segment independently without inter-arc coupling.
$\mathcal{L}_{\text{side}}$ (row~4) is the primary corrector of bilateral asymmetry, reducing
it from 2.259 to 0.769\,mm. Adding $\mathcal{L}_{\text{center}}$ alone (row~5) causes a
temporary regression in left arc (2.859~$\to$~4.049\,mm) because the center-alignment gradient
can reinforce incorrect arc assignments when bilateral disambiguation is still partial.
$\mathcal{L}_{\text{contrast}}$ (row~6) resolves this by establishing an unambiguous
separation between predicted arc midpoints; after which the center-alignment signal pulls each
prediction toward its correct target. The full combination (row~6) achieves the best result
on every metric simultaneously, with bilateral arc asymmetry reduced to 0.350\,mm.

Crucially, all six validation loss curves converge to a narrow band
(0.001028--0.001080)---a spread of only $5.2{\times}10^{-5}$---while the sum of three errors
ranges from 9.321 to 13.493\,mm, a 44.7\% relative difference invisible to the heatmap
validation loss. This demonstrates that heatmap validation MSE is an unreliable discriminator
among these configurations, and motivates reporting clinical measurement error as the primary
evaluation metric.

\subsection{Synthetic Data Ablation}
\label{sec:synthetic}

To quantify the contribution of VLM-generated synthetic images, we trained the full ACDentNet
model with and without the 20 synthetic training samples under identical 5-fold cross-validation
conditions. The synthetic images included an explicit 5\,mm calibration scale marker positioned
consistently with the clinical acquisition protocol. Table~\ref{tab:synthetic} reports results.

%% -------------------------------------------------------
%% TABLE 4 — Real+synthetic row updated to reflect thesis means.
%% *** ACTION REQUIRED: Run training without synthetic data and fill in real-only row. ***
%% -------------------------------------------------------
\begin{table}[htbp]
\caption{Effect of Synthetic Data Augmentation (ResNet-50, 5-Fold CV Mean, mm). Lower is
better. $^*$Run required: train without synthetic data and replace placeholder values.}
\begin{center}
\begin{tabular}{|l|c|c|c|}
\hline
\textbf{Training data} & \textbf{Incisor} & \textbf{L Arc} & \textbf{R Arc} \\
\hline
Real only (no synthetic)$^*$   & -- & -- & -- \\
Real + synthetic (VAE-aligned) & \textbf{2.347} & \textbf{2.324} & \textbf{2.690} \\
\hline
\end{tabular}
\label{tab:synthetic}
\end{center}
\end{table}

\subsection{Backbone Ablation}

% FIX C2/S8: Evaluation protocol for backbone table stated explicitly.
Table~\ref{tab:backbone_results} presents backbone ablation within the proposed framework.
These results were obtained using the full composite objective on a fixed 80/20 training/test
split (fold~1 partition) rather than 5-fold cross-validation, which accounts for the lower
absolute errors compared to the 5-fold means in Table~\ref{tab:cv_results}. ResNet-50 achieves
the most consistent performance across all five clinically relevant metrics. Although
ConvNeXt-Tiny achieves the lowest left arc error (0.84\,mm, underlined), it shows
substantially higher right arc error (2.74\,mm) and mandibular incisor error (2.26\,mm),
indicating bilateral inconsistency; ResNet-50 was therefore selected as the primary backbone.
HRNet exhibited unstable convergence under the current dataset size.

%% TABLE 5 — Backbone table unchanged; not modified by thesis.
\begin{table}[htbp]
\caption{Backbone Ablation (Full Composite Objective, Fold~1 Split, Mean Error in mm). Lower
is better. $\dagger$: ConvNeXt-Tiny achieves the lowest L Arc but shows bilateral
inconsistency (high R Arc and Incisor); ResNet-50 achieves uniformly low error across all
metrics. Note: these errors are lower than Table~\ref{tab:cv_results} because they reflect a
single fold rather than a 5-fold mean.}
\begin{center}
\begin{tabular}{|l|c|c|c|c|c|}
\hline
\textbf{Backbone} & \textbf{Mand.} & \textbf{L} & \textbf{R} & \textbf{Max} & \textbf{Max} \\
 & \textbf{Inc.} & \textbf{Arc} & \textbf{Arc} & \textbf{L} & \textbf{R} \\
\hline
ConvNeXt-Tiny$^\dagger$ & 2.26 & \underline{0.84} & 2.74 & 0.85 & 0.96 \\
EfficientNet-B3         & 2.30 & 2.79 & 5.92 & 0.71 & 1.75 \\
HRNet                   & 86.86 & 13.01 & 10.17 & 7.50 & 12.14 \\
ResNet-101              & 2.08 & 2.24 & 3.92 & 0.59 & 5.18 \\
ResNet-18               & 2.92 & 4.13 & 8.10 & 2.11 & 2.41 \\
\textbf{ResNet-50}      & \textbf{1.73} & \textbf{1.13} & \textbf{1.21} & \textbf{0.55} & \textbf{0.99} \\
\hline
\end{tabular}
\label{tab:backbone_results}
\end{center}
\end{table}

\subsection{Failure Analysis}

Failure cases were primarily observed in images with severe crowding or indistinct molar
boundaries, where visual ambiguity increases localization variance. In these cases arc
endpoint predictions shift anteriorly, elongating predicted arc length relative to ground
truth. The bilateral geometry constraints remain effective in all failure cases: arc ordering
is preserved and errors remain localized to arc measurements without propagating to incisor
or scale estimates. A second failure pattern involves maxillary images with minimal soft
tissue contrast, where predicted arc midpoints drift toward the gingival margin; this affects
both arcs symmetrically so the bilateral constraints neither exacerbate nor resolve it.

In additional experiments using a 60--20--20 (train/validation/test) split, measurement
errors increased by approximately 0.6\,mm compared to the primary cross-validation setup,
consistent with a reduction in effective training set size.

\subsection{Clinical Relevance}

The proposed method achieves a mean incisor error of 2.347\,mm and arc errors of 2.324\,mm
(left) and 2.690\,mm (right) across five folds, with best-fold errors reaching 2.024\,mm
(incisor), 2.066\,mm (left arc), and 2.045\,mm (right arc). Automatic tooth segmentation on
digital dental models has demonstrated high intra-rater reliability with ICC values of
0.987--0.997 \cite{im2022toothsegmentation}. Inter-observer variability in manual cast-based
orthodontic measurement is reported at 0.5--1\,mm \cite{schwendicke2020artificial}; the
current system's errors exceed that range, indicating further improvement is needed before the
system can serve as a direct replacement for cast measurement. However, in resource-limited
primary care settings where physical models and intraoral scanners are unavailable, the
achieved 2--3\,mm mean error provides clinically actionable information where none currently
exists. The approach requires no radiographic equipment, no casts, and no laboratory processing
time; it operates on photographs routinely captured at clinical visits with the addition of a
5\,mm calibration scale.

%% ==========================================================
\section{Conclusion}
%% ==========================================================

% FIX M7: Conclusion no longer overclaims "each component individually contributes."
We presented ACDentNet (Anatomy-Constrained Dental Network), a landmark learning framework for
automated mixed dentition space analysis from RGB intraoral photographs. By combining
heatmap-based localization with a six-component geometry-aware training objective, the method
enforces bilateral arc ordering and improves structural stability under limited-data conditions.
Per-term ablation on fold~1 demonstrates that the full combination of all six loss terms
achieves the best result across every metric simultaneously; the three bilateral constraint
terms function as a coordinated mechanism in which each term's contribution depends on the
others. Five-fold cross-validation yields mean errors of 2.347\,mm (incisor), 2.324\,mm (left
arc), and 2.690\,mm (right arc), demonstrating consistent measurement accuracy across
maxillary and mandibular arches. Operating solely on routine 2D photographs, the approach
provides a cost-effective alternative to cast- or scan-based workflows suitable for
low-resource clinical settings. Future work will focus on expanding the dataset across diverse
sites and on prospective validation before clinical deployment.

\section*{Acknowledgment}

The authors acknowledge the institutional ethics board approval and the participating patients
and guardians for their consent.

%% ==========================================================
\begin{thebibliography}{00}
%% ==========================================================

\bibitem{moyers1988handbook}
R.~E. Moyers, \textit{Handbook of Orthodontics}, 4th~ed.
Chicago: Year Book Medical Publishers, 1988.

\bibitem{tanaka1974prediction}
M.~M. Tanaka and L.~E. Johnston,
``The prediction of the size of unerupted canines and premolars in a contemporary orthodontic
population,''
\textit{Journal of the American Dental Association}, vol.~88, no.~4, pp.~798--801, 1974.

\bibitem{he2016deep}
K.~He, X.~Zhang, S.~Ren, and J.~Sun,
``Deep residual learning for image recognition,''
in \textit{Proc.\ IEEE CVPR}, 2016, pp.~770--778.

\bibitem{newell2016stacked}
A.~Newell, K.~Yang, and J.~Deng,
``Stacked hourglass networks for human pose estimation,''
in \textit{Proc.\ ECCV}, 2016, pp.~483--499.

\bibitem{sun2018integral}
X.~Sun, B.~Xiao, F.~Wei, S.~Liang, and Y.~Wei,
``Integral human pose regression,''
in \textit{Proc.\ ECCV}, 2018, pp.~529--545.

\bibitem{li2022simcc}
Y.~Li, S.~Yang, P.~Liu, S.~Zhang, Y.~Wang, Z.~Wang, W.~Yang, and S.-T.~Xia,
``SimCC: A simple coordinate classification perspective for human pose estimation,''
in \textit{Proc.\ ECCV}, 2022, pp.~89--106.

\bibitem{maji2022yolo}
D.~Maji, S.~Nagori, M.~Mathew, and D.~Poddar,
``YOLO-Pose: Enhancing YOLO for multi-person pose estimation using object keypoint similarity
loss,''
in \textit{Proc.\ IEEE/CVF CVPR Workshops}, 2022, pp.~2637--2646.

\bibitem{payer2016regressing}
C.~Payer, D.~\v{S}tern, H.~Bischof, and M.~Urschler,
``Regressing heatmaps for multiple landmark localization using CNNs,''
in \textit{Proc.\ MICCAI}, 2016, pp.~230--238.

\bibitem{wang2016benchmark}
C.~W. Wang, C.~T. Huang, J.~H. Lee, C.~H. Li, S.~W. Chang, M.~J. Siao, T.~M. Lai,
B.~Ibragimov, T.~Vrtovec, O.~Ronneberger, P.~Fischer, R.~Cootes, and C.~Lindner,
``A benchmark for comparison of dental radiography analysis algorithms,''
\textit{Medical Image Analysis}, vol.~31, pp.~63--76, 2016.

\bibitem{oktay2018anatomically}
O.~Oktay, E.~Ferrante, K.~Kamnitsas, M.~Heinrich, W.~Bai, J.~Caballero, S.~A. Cook,
A.~de Marvao, T.~Dawes, D.~P. O'Regan, B.~Kainz, B.~Glocker, and D.~Rueckert,
``Anatomically constrained neural networks (ACNN): Application to cardiac image enhancement
and segmentation,''
\textit{IEEE Trans.\ Medical Imaging}, vol.~37, no.~2, pp.~384--395, 2018.

\bibitem{li2020graphlandmark}
W.~Li, Y.~Lu, J.~Zheng, H.~Liao, C.~Lin, W.~Luo, J.~Shi, C.~Cheng, A.~Yuille,
L.~Zhou, and D.~Xu,
``Structured landmark detection via topology-adapting deep graph learning,''
in \textit{Proc.\ ECCV}, vol.~12354, 2020, pp.~266--283.

\bibitem{schwendicke2020artificial}
F.~Schwendicke, W.~Samek, and J.~Krois,
``Artificial intelligence in dentistry: Chances and challenges,''
\textit{Journal of Dental Research}, vol.~99, no.~7, pp.~769--774, 2020.

\bibitem{schwendicke2019deep}
F.~Schwendicke, T.~Golla, M.~Dreher, and J.~Krois,
``Convolutional neural networks for dental image diagnostics: A scoping review,''
\textit{Journal of Dentistry}, vol.~91, p.~103226, 2019.

\bibitem{im2022toothsegmentation}
J.~Im, J.-Y. Kim, H.-S. Yu, K.-J. Lee, S.-H. Choi, J.-H. Kim, H.-K. Ahn, and J.-Y. Cha,
``Accuracy and efficiency of automatic tooth segmentation in digital dental models using deep
learning,''
\textit{Scientific Reports}, vol.~12, no.~1, p.~9429, 2022.

\bibitem{ghorbani2025automated}
Z.~Ghorbani et al.,
``A novel deep learning-based model for automated tooth detection and numbering in mixed and
permanent dentition in occlusal photographs,''
\textit{BMC Oral Health}, vol.~25, p.~455, 2025.

\bibitem{kartbak2025classification}
S.~B.~A. Kartbak et al.,
``Classification of intraoral photographs with deep learning algorithms trained according to
cephalometric measurements,''
\textit{Diagnostics}, vol.~15, no.~9, p.~1059, 2025.

\bibitem{nguyen2025segmentanytooth}
K.~D. Nguyen et al.,
``SegmentAnyTooth: An open-source deep learning framework for tooth enumeration and
segmentation in intraoral photos,''
\textit{Journal of Dental Sciences}, vol.~20, no.~2, pp.~1110--1117, 2025.

\bibitem{koch2026angleclassification}
P.~J. Koch et al.,
``Deep learning for Angle classification based on intraoral photographs: An interpretability
perspective,''
\textit{BMC Oral Health}, vol.~26, p.~210, 2026.

\bibitem{kingma2014vae}
D.~P. Kingma and M.~Welling,
``Auto-encoding variational Bayes,''
in \textit{Proc.\ ICLR}, 2014.

\bibitem{cootes1995asm}
T.~F. Cootes, C.~J. Taylor, D.~H. Cooper, and J.~Graham,
``Active shape models---their training and application,''
\textit{Computer Vision and Image Understanding}, vol.~61, no.~1, pp.~38--59, 1995.

\bibitem{brody2022gatv2}
S.~Brody, U.~Alon, and E.~Yahav,
``How attentive are graph attention networks?,''
in \textit{Proc.\ International Conference on Learning Representations (ICLR)}, 2022.

\end{thebibliography}

\end{document}